\section{Conclusion}
\label{sec:conclusion}

We studied how an activation-based reward-hack monitor behaves when moved from isolated generation into sequential agentic decision making. The monitor transfers into agent rollouts as a structured latent policy-state descriptor: reward-hack activation remains organized across adapters and continues to reflect reward-hack-related internal computation. However, action-level risk is not a direct readout of this signal. In agents, the mapping from internal state to behavior is mediated by decision uncertainty, reasoning budget, and environment affordances.

Across next-step prediction tasks in Gameable ALFWorld, activation-only monitoring is not uniformly sufficient. Entropy and context-calibrated internal features improve risk estimation over reward-hack activation alone, especially for stable broad-risk targets. Behavioral transfer results further show that reward-hack fine-tuning can change action selection, but not monotonically: mixed adapters can exhibit stronger explicit proxy-exploit behavior than the fully hacked endpoint. Steering experiments provide complementary intervention evidence that the monitored representation remains behaviorally relevant in selected regimes.

Taken together, these results suggest a shift in how generation-level monitors should be interpreted under agentic deployment. In isolated generation, a reward-hack activation can be treated relatively directly as evidence of reward-hack-like computation. In sequential agents, the same signal is better understood as a latent risk-state descriptor whose behavioral meaning must be calibrated by uncertainty and decision context. Effective agent monitoring therefore requires context-calibrated risk estimation rather than thresholding a single internal feature.
